% 来源：OEBPS/Text/01_20.xhtml（OPF spine 顺序；仅第一册）
\cleardoublepage
\BookEntry[史蒂芬·平克《语言本能》]{史蒂芬·平克《语言本能》}

\BookCover{images/书单06语言本能.jpg}{[加] 史蒂芬·平克著| 欧阳明亮译| 浙江人民出版社2015年版}

\begin{BookSummaryBox}
\textbf{第一本真正讲述人类语言进化奥秘的书}
\par
\end{BookSummaryBox}

\subsubsection*{这本书谈什么？}

史蒂芬·平克以其深厚的语言学、认知神经学和进化心理学知识，证明了这样一个事实，即“语言能力是人类的本能”。他知道一般人对于语言学问题的困惑何在，真实地说明了我们只要张开嘴、发出声音，就可以把心中的想法传达到对方心中，而我们常常忘了这种能力其实是个奇迹。语言的复杂性常常让语言学家困扰，这也正是源于语言是人类本能的第一个原因。即便是没有口头语言的人，还可以有思想语言，因而习得语言就是知道如何将思想语言翻译成一串字词，或将字词转换成思想。每个人头脑中都有字和概念的组合，这套规则可以把字组合起来，并将概念之间的关系传递出去。
\par

平克用令人信服且生动有趣的例证，展示了一场常识对谬论的彻底胜利。
\par

\subsubsection*{芒格为什么谈起这本书？}

在“论基本的、普世的智慧”演讲中，芒格说到有件让他获益匪浅的事情：
\par

\begin{quotation}
一个在香港一所名牌中学当校长的朋友送了我一本书，叫《语言本能》，作者是史蒂芬·平克。平克是一个语义学教授，他的名气没有诺姆·乔姆斯基那么大。
\par

平克说，人类的语言能力不仅仅是后天学来的——从很大程度上来说，它还跟先天的遗传有关。其他动物，包括黑猩猩，都缺乏真正有用的语言基因。语言是上天赐给人类的礼物。平克很漂亮地证明了他的观点。
\par

当然，乔姆斯基也已经证明这一点。只有非常愚蠢的人才不明白语言能力大部分来自人类基因的道理。平克无法理解为什么乔姆斯基这样的天才居然还认为语言能力是否基于人类的基因尚无定论。实际上，平克是这么说的：“什么尚无定论，活见鬼了！人类得到语言本能的途径跟得到其他本能的途径完全一样——那就是达尔文的自然选择。”
\par
\end{quotation}

在芒格看来，因为乔姆斯基脑中强烈的意识形态导致了他不愿意承认这一事实。
\par

\subsubsection*{《语言本能》作者史蒂芬·平克如是说}

“人类创造语言用于诉说自己内心的不满。”
\par

“婴儿是天生的‘语法专家’。”
\par

“我们本以为语言是文化的产物，错误地认为教育程度偏低的人群语言表达能力也会差一些。但实际上，美国黑人在英语表达方面却更胜一筹。”
\par

“人们很容易高估语言的能力，以为语言决定着我们的思维。实际上，语言不是思维的唯一方式。无论是英语还是其他任何自然语言，都不能用作心智计算的介质。心语，才是思维的语言。”
\par

“一些研究显示，给住院的病人送花不仅是一种温暖的问候，它的确能改善患者的心情，缩短康复的时间。”
\par

“父母对孩子最重要的影响在受孕的那一刻就已经完成，孩子是在和同龄人的交往互动中习得文化、发展个性的。”
\par

